\section{Methods}

\subsection*{Datasets}

GSE134839 \citep{aissa2021}: PC9 cells treated with
\SI{2}{\micro\molar} erlotinib; six Drop-seq DGE matrices
(D0, D1, D2, D4, D9, D11); 6{,}508 total cells.
GSE150949 \citep{oren2021}: PC9 cells with Watermelon lineage tracing
under osimertinib; 56{,}419 cells across D0/D3/D7/D14; subsampled to
5{,}000.
GSE243562 \citep{wurtz2024}: two EGFR-mutant PDX models (YU-003,
YU-006) under osimertinib; 22{,}032 cells across four samples
(untreated and residual disease for each model).
GSE131907 \citep{kim2020}: treatment-naive patient LUAD atlas;
208{,}506 cells across 44 patients; subsampled to 5{,}000 tLung
epithelial cells, then a further 1{,}500 cells for tractable
Vietoris--Rips computation (filter and seed details in
Supplementary Methods~S1).

\subsection*{Preprocessing}

We used scanpy 1.11 \citep{wolf2018} with standard QC
(\texttt{min\_genes} $= 200$, \texttt{max\_genes} $= 5000$,
\texttt{max\_pct\_mito} $= 20\%$, \texttt{min\_cells} $= 3$),
total-count normalisation to 10{,}000 counts per cell, log1p
transformation, 3{,}000 highly-variable genes, regression of total
counts and mitochondrial percentage, scaling to max 10, and PCA with
50 components (\texttt{random\_state} $= 42$). HVG batch keys per
cohort: \texttt{timepoint} for GSE134839 and GSE150949,
\texttt{pdx} for GSE243562, and \texttt{Sample} for GSE131907
(rationale in Supplementary Methods~S2). For PDX data, mouse genes were removed
by a symbol-pattern heuristic (Supplementary Methods~S3; 29\% of
features filtered; limitations in Discussion). Cell-cycle scores were
computed from \citet{tirosh2016} S- and G2/M-phase gene lists via
\texttt{sc.tl.score\_genes\_cell\_cycle}.

\subsection*{Persistent homology}

We used \texttt{ripser} 0.6 \citep{bauer2021,tralie2018} to compute the
Vietoris--Rips filtration up to $H_1$ with the Euclidean metric on the
top-30 PCA coordinates; homology is computed over $\mathbb{F}_2$
coefficients (ripser default). Maximum persistence is defined as
$\max(\text{death} - \text{birth})$ over finite $H_1$ bars; this is a
hybrid statistic whose existence is a topological invariant but whose
numerical scale is metric-dependent (Supplementary Methods~S4).
Wasserstein distances between diagrams were computed with
\texttt{persim} 0.3 using the 2-Wasserstein metric. The $p$-Wasserstein
distance between persistence diagrams is stable under bounded
perturbations of the input point cloud given a finite-degree-$p$
total-persistence hypothesis \citep{cohen-steiner2007,cohen-steiner2010}.

\subsection*{Principal component count sensitivity}

Max $H_1$ depends on the number of PCs retained: across
$n_{\text{PCs}} \in \{10, 20, 30, 50\}$ absolute values vary within a
factor of two, but the drug-treated $>$ untreated ordering is
preserved at every PC count. The $n_{\text{PCs}} = 30$ default
follows prior scTDA work \citep{rizvi2017}. Full sweep numbers and
the rationale for the hybrid-statistic framing are in Supplementary
Methods~S4--S5 and Supplementary Table S5.

\subsection*{Permutation tests}

For each group, we shuffled expression values within each cell
(permuting gene labels independently for every cell), re-fit PCA using
\texttt{sklearn}, and recomputed $H_1$ max persistence. The $p$-value
is the fraction of null permutations with max $H_1 \geq$ the observed
value. Discovery cohort used 500 permutations per timepoint; validation
cohorts used 100 permutations on 1{,}000-cell subsamples for
tractability.

\subsection*{Mapper}

We used \texttt{kmapper} 2.0 \citep{vanveen2019} with an auto-tuned DBSCAN clusterer
(\texttt{eps} set to the 50th percentile of 10th-nearest-neighbor
distances in each cover element). Default parameters were
$n_{\text{cubes}} = 15$ and $\text{perc\_overlap} = 0.3$. Parameter
sensitivity was assessed across
$n_{\text{cubes}} \in \{10, 15, 20, 25, 30\}$ and overlap $\in \{0.2,
0.3, 0.4, 0.5\}$.

EMT score (used as Mapper filter for the cell-line analysis) was
computed via \texttt{scanpy.tl.score\_genes} on a 33-gene MSigDB
Hallmark EMT cassette (28 of 33 present in GSE134839; full gene list
in Supplementary Methods~S6).

\subsection*{Gene connectivity}

For each gene $g$ and Mapper node $N$ with cells $C_N$, we computed
$|\text{mean}(g \mid C_N) - \text{mean}(g \mid \text{all cells})|
\times |C_N|$, summed over all nodes, and divided by total cells.
Higher scores indicate genes that vary more between Mapper nodes than
globally.

\subsection*{Pathway enrichment}

Gene set enrichment was performed with \texttt{gseapy 1.1} against
MSigDB\_Hallmark\_2020 and GO\_Biological\_Process\_2021.

\subsection*{Computational environment}

All analyses ran on Python 3.11 in a conda environment specified in
\texttt{envs/environment.yml}; random seed 42 is used for every
stochastic operation. Full package versions and platform notes
(macOS arm64 and Linux x86\_64): Supplementary Methods~S14.

\subsection*{Pipeline architecture and unit tests}

The framework is organised as fifteen numbered modules
(\texttt{01\_preprocess.py} through \texttt{15\_maynard\_full\_cohort.py})
that each accept a standard \texttt{AnnData} object and produce typed
intermediate outputs. Users can substitute their own datasets by
providing an \texttt{AnnData} file with a \texttt{timepoint}
observation column. Full module list and the \texttt{pytest} unit-test
suite (8 tests, all passing) are described in Supplementary
Methods~S8--S9.

\subsection*{Code availability}

Code is available at \url{https://github.com/htlin222/sctda-cancer-plasticity} (MIT license). The version-tagged software release accompanying this submission is \texttt{v2.7.1-submission}, archived as a public GitHub release at \url{https://github.com/htlin222/sctda-cancer-plasticity/releases/tag/v2.7.1-submission}. The complete pipeline is reproducible via \texttt{make pipeline} from raw GEO data to all figures and tables. A unified \texttt{sctda} CLI is planned for v2.0; the current pipeline is invoked via the Makefile and individual \texttt{scripts/0X\_*.py} files. A \texttt{CITATION.cff} file at the repository root provides machine-readable citation metadata for software-citation tooling.

\subsection*{Data availability}

All raw scRNA-seq data used in this study are publicly available from
the Gene Expression Omnibus under accessions GSE134839, GSE150949,
GSE243562, and GSE131907; the Maynard et al.\ 2020 EGFR-mutant cohort
was obtained from the authors' public Google Drive archive referenced in
\citet{maynard2020}. Processed AnnData objects, persistence diagrams,
and Mapper graphs are regenerable end-to-end from the raw GEO data by
running \texttt{make pipeline}; seeds and parameters are pinned in the
repository so derived outputs are bit-stable across runs.
